\chapter{Implementation Details}

\section{Frontend Implementation}

\subsection{Component Architecture}

The React frontend follows a feature-based folder structure:

\begin{lstlisting}[language=bash, caption={Frontend Folder Structure}]
frontend/src/app/
  components/
    interview/    # VideoRecorder, SpeechRecognition, AIAvatar
    ui/           # Radix UI wrappers (Button, Card, Dialog, ...)
  hooks/          # useSocket, useAudioAnalysis, useWebRTC
  pages/          # One file per route
  services/       # api.ts, interview.ts, resume.ts
  stores/         # authStore.ts, interviewStore.ts
  types/          # TypeScript interfaces
\end{lstlisting}

\subsection{State Management}

Zustand stores manage global application state. The \texttt{interviewStore} is the most complex, managing:
\begin{itemize}
    \item Current interview session and question.
    \item Recording state and media stream.
    \item Loading and error states.
    \item Actions: \texttt{createInterview}, \texttt{startInterview}, \texttt{getNextQuestion}, \texttt{submitResponse}, \texttt{endInterview}.
\end{itemize}

A key design decision was to \textbf{not} call \texttt{getNextQuestion} automatically after \texttt{submitResponse}. This prevents duplicate API calls when both the store and the page component try to fetch the next question simultaneously.

\subsection{VideoRecorder Component}

The \texttt{VideoRecorder} component handles camera initialisation with a three-attempt fallback strategy:
\begin{enumerate}
    \item Ideal quality: 1280×720, echo cancellation, noise suppression.
    \item Basic: \texttt{\{video: true, audio: true\}}.
    \item Audio only: \texttt{\{video: false, audio: true\}}.
\end{enumerate}

The \texttt{facingMode: 'user'} constraint was removed because it causes black screens on desktop GPUs with integrated + discrete GPU configurations. The component waits for the \texttt{loadedmetadata} event before calling \texttt{play()} to ensure the first frame is visible.

\subsection{SpeechRecognition Component}

The \texttt{SpeechRecognition} component uses the native Web Speech API with:
\begin{itemize}
    \item \texttt{continuous: true} and \texttt{interimResults: true} for real-time transcription.
    \item Automatic restart on \texttt{no-speech} events (silence does not stop recognition).
    \item Callback refs (\texttt{useRef}) for all event handlers to prevent stale closure issues.
    \item CSS-only animated waveform bars (no \texttt{Math.random()} to avoid infinite re-renders).
\end{itemize}

\section{Backend Implementation}

\subsection{Interview Route Architecture}

The interview route implements a \textbf{fire-and-return} pattern for response submission:

\begin{lstlisting}[language=JavaScript, caption={Fire-and-Return Response Submission}]
// Step 1: Save response immediately
await Interview.findByIdAndUpdate(id, {
  $push: { responses: responseData }
});

// Step 2: Return 200 instantly
res.json({ success: true, message: 'Response submitted' });

// Step 3: Run AI analysis in background (non-blocking)
(async () => {
  const analysis = await geminiService.analyzeResponse({...});
  await Interview.findByIdAndUpdate(id, {
    $set: { 'analysis.contentMetrics.relevanceScore': analysis.scores.relevance, ... }
  });
})();
\end{lstlisting}

\subsection{Feedback Generation with Synchronous Fallback}

The feedback route detects if background analysis was incomplete and runs synchronous analysis:

\begin{lstlisting}[language=JavaScript, caption={Synchronous Analysis Fallback}]
const hasScores = metrics && (
  (metrics.relevanceScore || 0) > 0 ||
  (metrics.technicalAccuracy || 0) > 0
);

if (!hasScores && interview.responses.length > 0) {
  // Run synchronous analysis on all responses
  for (const resp of interview.responses) {
    const analysis = await geminiService.analyzeResponse({...});
    allScores.push(analysis.scores);
  }
  // Average scores and update DB
  await Interview.findByIdAndUpdate(id, { $set: { ... } });
}
\end{lstlisting}

\subsection{Code Execution Service}

The code execution service supports two backends:
\begin{itemize}
    \item \textbf{Piston API} (default): Free public API, no authentication required.
    \item \textbf{Judge0} (optional): Requires API key, more reliable for production.
\end{itemize}

The Python test harness uses a four-attempt fallback chain:
\begin{enumerate}
    \item \texttt{fn(*expanded\_args)} -- normal call with expanded arguments.
    \item \texttt{fn(*[expanded\_args])} -- wrapped in list.
    \item \texttt{fn(*original\_args)} -- raw original arguments.
    \item \texttt{fn(*[original\_args])} -- raw wrapped.
\end{enumerate}

\section{AI Server Implementation}

\subsection{Eye Contact Scoring}

Eye contact is scored using MediaPipe Face Mesh iris landmarks (indices 468 and 473):

\begin{lstlisting}[language=Python, caption={Eye Contact Scoring}]
def _estimate_eye_contact(self, landmarks) -> float:
    left_iris = landmarks.landmark[468]
    right_iris = landmarks.landmark[473]
    nose_tip = landmarks.landmark[1]
    
    # Horizontal deviation from nose (camera center proxy)
    avg_dev = (abs(left_iris.x - nose_tip.x) + 
               abs(right_iris.x - nose_tip.x)) / 2.0
    avg_vdev = (abs(left_iris.y - nose_tip.y) + 
                abs(right_iris.y - nose_tip.y)) / 2.0
    
    h_score = max(0.0, 1.0 - (avg_dev / 0.12)) * 100
    v_score = max(0.0, 1.0 - (avg_vdev / 0.12)) * 100
    return round((h_score * 0.6 + v_score * 0.4), 1)
\end{lstlisting}

\subsection{Posture Scoring}

Posture is scored using MediaPipe Pose landmarks for shoulder alignment, spine offset, hip level, and forward lean:

\begin{equation}
\text{Posture Score} = 0.3 \times S_{shoulder} + 0.3 \times S_{spine} + 0.2 \times S_{hip} + 0.2 \times S_{lean}
\end{equation}

where each sub-score is computed as $\max(0, 1 - \frac{\text{deviation}}{\text{threshold}}) \times 100$.
